\chapter{厨房、药箱和花园里都藏着化学}

\begin{kaichang}
不用去实验室，今天只做三件事，都在厨房里。

第一件：找一只透明的玻璃杯，倒半杯温水，加一勺白糖，搅一搅。糖不见了——水还是清澈透明的，但尝一口，是甜的。

第二件：找一口干净的小锅，放一勺白糖，小火慢慢加热（这一步请家长来做，或者看视频，熔化的糖温度很高，沾到皮肤会严重烫伤）。几分钟后，白色的糖先熔化成透亮的液体，再变成金黄，再变成深棕，飘出一股焦香。尝一点放凉了的焦糖，又香又苦，跟白糖完全不是一个味儿。

第三件：一杯温糖水，撒一小撮做馒头用的干酵母，轻轻搅一下，放在一边。半小时后回来，水面上浮起一层泡沫，杯壁挂满小气泡，像一杯慢慢煮开的水——但它是凉的。凑近闻，有一股淡淡的酒味。

同一勺糖，三种遭遇。请你猜一猜：这三件事里，哪些只是糖“变了样子”，哪些是糖真的“变成了别的东西”？酵母又凭什么让糖水冒泡——它是活的，一个活物做的事，跟化学有什么关系？把你现在的猜测写在纸上，别偷看后面的答案。这一章结束时，我们会带着一副全新的眼镜回来检查它。
\end{kaichang}

\section{先只说看得见的}

科学的起点不是术语，而是诚实的观察。所以在引入任何概念之前，我们先把三件事里“能看见、能闻到、能尝到、能测到”的部分老老实实记下来，一个字都不解释。

\textbf{糖水。}糖放进水里，颗粒越变越小，最后彻底看不见了。水仍然是透明的，颜色没变，没有气泡，摸一摸杯壁，温度几乎没有变化。可是糖水明明是甜的，糖到哪里去了？如果你有足够的耐心，把糖水放在通风处让水慢慢蒸发，几天后杯底会重新留下白色的固体，尝一尝——还是糖。糖好像只是“藏”进了水里，随时可以被找回来。

\textbf{焦糖。}加热后的糖讲的是另一个故事。颜色从白到金黄再到棕黑，气味从无味到焦香，放凉后得到的固体掰开来，质地、味道、气味全都变了。更麻烦的是：你把焦糖放进水里，它倒是能溶掉一部分，可无论你怎么蒸发、怎么冷却，都再也得不到原来那种白糖了。加热时锅上方还飘走了水汽，熬过头了甚至会冒出呛人的烟。原来的糖似乎真的消失了，换成了一批新面孔。

\textbf{酵母糖水。}这件事里最抢眼的迹象是气泡：有气体在杯子里源源不断地生成。这气体是哪来的？同样一杯糖水，不放酵母，摆在旁边半天，什么也不会发生——所以气泡既不是水自己变出来的，也不是原来就躲在杯子里的。再看别的迹象：杯子摸上去微微温热；甜味在一天天变淡；放上一两天，能尝到淡淡的酒味。糖在减少，某种新东西在增加，还有一个看不见摸不着的“劳工”——酵母——在加班。

观察记录到此为止。注意一件事：这三段描述里我没有用任何化学名词，但你多半已经感觉到，这三件事“性质不太一样”——第一件好像能倒回去，后两件倒不回去。请珍惜这个直觉，它非常值钱。接下来整章的任务，就是给这种“一样还是不一样”找到一个站得住的判据。

\section{把镜头推进到粒子}

只盯着杯子看，永远也看不出“能不能倒回去”的道理。得把镜头往里推，推到比头发丝细一万倍、再细一万倍的尺度。

这本书从第 6 章起会正式讲述原子的故事——它是怎么被一步步逼出来的，证据是什么。现在你只需要临时接受一幅图像，把它当成一张先赊账、后兑现的地图：

\begin{itemize}[itemsep=2pt]
\item 世界上一切物质都由极其微小的粒子组成；
\item 糖、水、酒精、二氧化碳这些“各有各的脾气”的物质，各自对应一种特定的\textbf{分子}——保持这种物质性质的最小单元；
\item 分子是由\textbf{原子}按固定方式搭起来的小结构。搭积木的块数少得可怜：糖分子里只有碳、氢、氧三种原子，但具体的搭法决定了它是糖而不是别的。
\end{itemize}

拿这幅图像重新看刚才的三件事。

\textbf{糖溶于水。}糖分子本身没有变。它们只是从紧紧挨在一起的糖粒里散开，均匀地混进了水分子之间的空隙里，像一盒积木被拆开、撒进一大盆玻璃珠——积木块一块都没坏。所以甜味还在，所以把水蒸干，糖分子重新聚拢，又变回白糖。

\textbf{糖加热变焦。}高温下，糖分子内部的搭法被打散了：一部分原子脱下“水”的成分，互相连接成更大、颜色更深的分子（焦香和棕黑色就来自它们），还有一部分变成了水汽和小分子飘走。这一回，积木不但被拆开，还被重新拼成了完全不同的造型。想让它们自己拼回“糖”？没有任何简单的办法。

\textbf{酵母发酵。}酵母是一种微生物——一个真正的、活着的细胞。它把糖分子“吞”进体内，在细胞里把糖分子拆开，重新组装成两类东西：二氧化碳分子（气泡的来源）和酒精分子（酒味的来源）。拆糖的过程释放出一些能量，供酵母生长繁殖，顺便把杯子焐得微微发热。请注意一个细节：气泡里的二氧化碳并不是酵母“无中生有”变出来的——组成它的碳原子和氧原子，原来就锁在糖分子里。酵母干的是拆解和重新组装的活，不是凭空创造原子。

到这里，三件事已经分成了清清楚楚的两堆：第一堆里，分子本身没有变；后两堆里，分子被拆开重组，出现了原来没有的新分子。这一道分界线，就是化学里最根本的分界线。

\section{物理变化与化学变化}

现在可以给这两堆起名字了。

\begin{itemize}[itemsep=2pt]
\item \textbf{物理变化}：物质的形状、状态、位置变了，但分子本身没变。糖溶于水、水结成冰、酒精挥发、把铁丝弯成圆圈——粒子还是那些粒子。
\item \textbf{化学变化}（也叫化学反应）：旧分子被拆散，原子重新组合，生成了新分子。糖变焦、糖发酵、铁生锈、木头燃烧、牛奶变酸——变化之后，世界上多出了原本不存在的分子种类。
\end{itemize}

判断一个变化属于哪一类，真正的判据只有一条：\textbf{有没有生成新物质（新分子）}。变化是否剧烈、是否发光发热、是否冒泡、能不能倒回去，全都只是旁证，不是判据本身。

用这个判据快速清点一下你身边的一天：早晨烧开水，壶嘴冒出的“白气”是水蒸气凝成的小水滴——水分子还是水分子，物理变化。煎鸡蛋，蛋清从透明变成白色固体——蛋白质分子的结构被高温永久改变，化学变化。切开的苹果放一会儿变黄变褐——果肉里的分子和空气中的氧发生了反应，化学变化。把盐撒进汤里——盐不见了但汤变咸，粒子只是分散开了，物理变化（这个例子后面还要回来找麻烦，它其实没那么简单）。晚上给手机充电——电池内部悄悄发生着化学变化，把电能存进分子的重新排布中。

你看，根本不用烧掉课本，也根本不用去实验室：你的厨房从早到晚都在运转一座化工厂。

\begin{shiyan}
\textbf{材料}：两个干净的矿泉水瓶（500 毫升左右）、温水（不烫手，约 $35\,^{\circ}\mathrm{C}$）、白糖、一小包做面食用的干酵母（超市有售）、两只气球、一支记号笔。

\textbf{步骤}：①两个瓶子各装半瓶温水，各加两勺白糖，摇匀。②给瓶子贴上标签：A 瓶加入半勺干酵母，摇一摇；B 瓶什么都不加。③在两只瓶口各套上一只气球，把瓶子放在温暖的地方（比如窗台边，别暴晒）。④每隔 15 分钟观察一次，持续一到两个小时，把看到的变化记下来。

\textbf{观察点}：哪只气球慢慢鼓起来了？B 瓶的气球呢？鼓起来的气球里装的是什么气体，它原来藏在哪里？再过一两天闻一闻 A 瓶（别喝），气味有什么变化？B 瓶为什么一直静悄悄的——它在这套装置里扮演什么角色？

\textbf{安全提示}：所有材料都可以入口，但实验后的液体敞口放置过，别喝。气球不要吹胀到极限，防止爆裂惊吓。水温不要超过 $40\,^{\circ}\mathrm{C}$——温度太高酵母会死，实验会“失败”，但这个失败本身值得想一想：活的东西为什么怕热？第 49 章讲酶的时候会回答你。

B 瓶是这个实验里最重要的设计：它和 A 瓶唯一的差别就是“没有酵母”，所以它替我们排除了“糖水自己会冒泡”“温水自己会冒泡”这些解释。这种设计叫\textbf{对照}，第 2 章会专门讲它为什么是让实验可信的关键。
\end{shiyan}

\section{化学家的三种语言}

同一件事，化学家可以用三种不同的语言来说，你得学会在它们之间翻译。

\textbf{第一种是宏观的语言}——眼睛、鼻子和手的语言：“白糖加热后变焦发黑，发出焦香。”这是我们一直在用的语言，它的优点是直接、诚实，缺点是说不出“为什么”。

\textbf{第二种是微观粒子的语言}：“糖分子在高温下拆解、脱水，原子重新组合成更大的深色分子。”这种语言能解释，但它谈论的东西谁也看不见——所以一整本书都在训练你“脑补”粒子画面的能力。

\textbf{第三种是符号的语言}，最省墨，也最唬人。每种分子可以写成一串符号：水写作 \chem{H_2O}（两个氢原子和一个氧原子），二氧化碳写作 \chem{CO_2}（一个碳原子和两个氧原子），葡萄糖写作 \chem{C_6H_{12}O_6}（6 个碳、12 个氢、6 个氧）。酵母发酵那件热闹的事，用符号写出来只有一行：

\[\chem{C_6H_{12}O_6} \rightarrow 2\chem{C_2H_5OH} + 2\chem{CO_2}\]

读法是这样的：箭头左边是“反应前”，右边是“反应后”；前面的数字 2 表示“两个分子”。整行念作：一个葡萄糖分子，被酵母拆解重组为两个酒精分子和两个二氧化碳分子。

现在做一件小孩子也能做、却很深刻的事：数一数箭头两边的原子。左边：6 个碳、12 个氢、6 个氧。右边：两个酒精分子共 4 个碳、12 个氢、2 个氧，两个二氧化碳分子共 2 个碳、4 个氧——加起来，6 个碳、12 个氢、6 个氧。\textbf{一个不多，一个不少。}化学变化不创造原子，也不消灭原子，它只是把原子重新分组。这行符号其实是一本账，怎样记账、怎样对账，第 26 章会专门教。

要提醒你的是：符号是地图，不是照片。\chem{H_2O} 这三个字符不告诉你水分子长什么形状、有多大、怎么动——那些要到第 20 章和第 23 章才揭晓。符号语言的力量恰恰在于它敢省略：它只写下最重要的信息——什么东西变成了什么东西，原子数目对不对得上。整本书里，我们将反复在宏观、微观、符号三种语言之间来回切换；当你发现自己能脱口而出一个现象的三套说法时，你就已经在用化学家的方式思考了。

\begin{qiaoliang}
粒子语言说糖分子“散开混进水分子之间”——到底是多少个？来估算一下。一小勺葡萄糖大约 5 克。第 5 章会正式介绍“摩尔”这个数数单位，这里先借用结论：180 克葡萄糖里约有 \ud{6.02\times 10^{23}}{} 个分子，那么 5 克里大约有

\[\frac{5}{180}\times 6.02\times 10^{23}\approx 1.7\times 10^{22}\ \text{个分子}\]

这个数字大到没有生活经验可以对照，硬比一下：把这一小勺糖里的分子平均分给全世界 80 亿人，每人能分到约 $2\times 10^{12}$ 个——两万亿个；它比地球上所有海滩和沙漠里的沙粒总数还要多。所以那杯糖水里并没有什么“糖的影子”或“糖的气味”，有的是一万亿亿个完好无损的糖分子，均匀地在水分子的缝隙里游荡。反过来你也该明白：既然每勺物质都裹着如此天文数字的粒子，化学家不可能一颗颗点名，必须靠符号和“摩尔”这样的批量计数法来工作——这就是第 5 章要解决的麻烦。
\end{qiaoliang}

\section{一碗焦糖和一位税务官}

\begin{zhengju}
“有没有新物质生成”这条判据听起来轻巧，可它背后站着一个硬问题：\textbf{你怎么知道生成了新东西？}颜色变了、冒泡了，都不算数——那什么才算数？18 世纪末，法国科学家拉瓦锡用一件毫不起眼的仪器给出了答案：天平。

当时的流行理论叫“燃素说”：能烧的东西里都含着一种叫“燃素”的物质，燃烧就是燃素跑掉的过程。这个说法能解释火焰，却有一个尴尬的细节——木头烧完变轻了（燃素跑了，合理），可铁、锡、铅烧过之后变重了（燃素跑了，东西应该变轻才对呀？）。燃素派的人只好硬着头皮说：燃素的重量是负的。

拉瓦锡的办法是把整个变化“关进账房”。他把锡放在密闭的玻璃容器里加热，先称总重量，加热后锡表面生成了灰白色的“锡灰”，锡明明发生了变化，可整个容器连同里面的空气一起称——总重量分毫不变。打开容器，空气嘶地涌进去，再称，容器变重了，增加的重量恰好等于锡灰比锡多出来的重量。结论只有一个：锡增加的那部分质量，来自空气。

更漂亮的实验在 1774 年前后：他把汞放在密闭容器里连续加热十二天，银亮的汞表面浮起一层红色粉末（当时叫“汞灰”），同时容器里的空气体积减少了约五分之一，剩下的气体能让蜡烛熄灭、让小鼠窒息。再把红色粉末收集起来加强热，它又变回汞，并放出一种气体——体积恰好等于之前“消失”的那五分之一，这种气体让烧红的木炭剧烈燃烧、让小鼠活蹦乱跳。拉瓦锡断言：空气不是铁板一块，燃烧也不是放出什么燃素，而是物质与空气中那五分之一的成分（他命名为氧）结合。

这个故事有三点值得记住。第一，拉瓦锡的“眼睛”是天平：他看不见分子，但他能通过称量追踪物质的去向——“颜色、气味、气泡只是线索”，而\textbf{精确的质量账目是证据}。第二，他没有靠一句雄辩推翻燃素说，而是给出了一组对方无法解释的数字；燃素说退出历史舞台，不是因为丢脸，是因为“负重量的燃素”在账目上再也圆不下去。第三，拉瓦锡本人也并不全对——他把“热”也列进了一张元素表，当作一种物质，这个错误几十年后才被纠正。科学不是一群从不出错的天才，而是一条不断对账、不断改错的长河。
\end{zhengju}

\section{这条分界线也有毛边}

讲了半天“物理变化还是化学变化”，现在该坦白：这条分界线很好用，但它不是一把没有毛边的刀。

\textbf{毛边一：判断本身需要手段。}“有没有新分子”听起来干脆，可分子看不见摸不着。糖溶于水，你凭什么敢说糖分子没变？得靠蒸发回收、靠称质量、靠更多仪器——下一章专门解决“怎样证明看不见的东西存在”。在证据不够的时候，诚实的回答是“现在判断不了”，而不是猜一个。

\textbf{毛边二：有些变化就骑在分界线上。}食盐溶于水：食盐是钠离子和氯离子整齐堆成的晶体，入水后晶体拆散，两种离子各自被一群水分子团团围住。粒子层面发生了相当可观的重新安排，可蒸干水，食盐原样回来。中学课本把它归为物理变化，你也该知道它归得有点勉强——第 25 章讲溶液时会重新翻开这笔账。

\textbf{毛边三：有些变化超出了这条线的管辖。}荧光灯通电发光，没有新分子生成，是物理过程；而放射性元素放出的射线里，原子的种类自己都变了——那已经不属于化学变化的范畴，是第 8 章要讲的原子核的故事。化学变化只管“原子怎么重新分组”，管不着“原子本身会不会变”。

\textbf{毛边四：可逆不可逆不是判据。}我们开场时靠“能不能倒回去”的直觉把三件事分了类，这个直觉常常灵验，但不是法律。给手机电池充电，放电时的化学变化被整个倒了回来；而冰化成水再冻成冰，照样是物理变化。能不能逆转，取决于能量和条件（第 27、28 章细讲），跟“是不是化学变化”是两回事。

所以请记住这条分界线的正确用法：它是侦探手里的放大镜，不是法官手里的锤子。先用它快速分类，遇到骑线案，就老老实实往粒子层和证据层再挖一层。

\begin{wuqu}
“看到冒气泡、变色、发光发热，就说明发生了化学变化。”——这是初中阶段最容易踩的坑。逐个反驳：烧开水剧烈冒泡，气泡里只是水蒸气，水分子没变，物理变化；拧开汽水瓶嘶嘶冒泡，那是原来溶解在水里的二氧化碳逃出来，分子没变，物理变化；白炽灯又亮又烫，是灯丝被电流烧热后发光，没有新物质，物理变化；红墨水把一杯水染红，颜色大变，分子还是原来那些分子。反方向同样成立：铁栏杆生锈，不冒泡、不发光、几乎不发热，慢得你毫无察觉，却是货真价实的化学变化。线索永远只是报案电话，不是判决书——想定案，请回到粒子层问一句：有没有新分子？
\end{wuqu}

\section{回到那三杯糖}

现在兑现开场的承诺，用新眼镜重新看那三件事，顺便批改你写下的猜测。

\textbf{糖水}：糖分子完好无损地分散进水分子之间。分子没变，是物理变化。甜味来自糖分子本身，所以水甜了；蒸发让水分子先走，糖分子重新聚成晶体，所以糖能找回来。

\textbf{焦糖}：糖分子在高温下脱水、拆解、重组，生成了几十种新的深色大分子和小分子。有新物质生成，是化学变化。颜色、焦香、苦味全都来自新分子；积木已经拼成了别的造型，所以白糖找不回来了。

\textbf{酵母糖水}：酵母细胞把葡萄糖分子拆解重组成酒精分子和二氧化碳分子，化学式写作 \chem{C_6H_{12}O_6} $\rightarrow$ $2\chem{C_2H_5OH}$ $+$ $2\chem{CO_2}$。有新物质生成，是化学变化。气泡是二氧化碳，酒味是酒精，温热是拆解过程中释放的能量。原子总数从头到尾没变——你随时可以按前面的方法数一遍账。

还有开场那个看似闲聊的问题：酵母是活的，它做的事跟化学有什么关系？现在可以认真回答了。生物学研究生命系统——细胞怎样生长、繁殖、感知环境；可酵母身体里的每一个动作，落到底层都是分子与分子的碰撞、拆解和重组。\textbf{生命过程不是化学的例外，而是化学最精巧的一种组织方式。}这也是为什么这本讲化学的书，最终会一路讲到细胞、DNA 和进化——那是第二册和第三册的故事。

\section{药箱、花园和一座看不见的工厂}

学到这里，你已经握住了侦探的第一件工具，是时候带着它出门巡逻了。

\textbf{药箱。}拿起任何一盒药，看成分表：比如常见的退烧止痛成分“对乙酰氨基酚”，化学式 \chem{C_8H_9NO_2}——它和白糖一样，是一种有确定分子结构的纯净物质，只不过它恰好能干预你身体里“疼痛信号”的化学反应。这里有两句话值得刻在心里。第一句：“化学成分”不等于“人工有害”——水、糖、维生素 C、你呼吸的氧，全都是化学成分，你自己就是一座由化学成分组成的、会走路的化工厂。第二句：“天然”不等于“安全”——毒蘑菇和发芽土豆里的毒素都是百分之百的天然产物，而救命的胰岛素是工厂里合成的。一种物质是好是坏，取决于它是什么、剂量多大、用在谁身上，而不是它姓“天然”还是姓“化学”。也正因为如此，药千万别按科普书的描述自行增减，剂量和个体差异是医生手里的专业问题。

\textbf{花园。}花盆里的土、肥料袋上的“氮磷钾”、叶片在阳光下进行的反应、堆肥桶里发热的腐殖质，全是化学。叶子把二氧化碳和水编织成糖——某种意义上是酵母发酵的“反向工程”，第 54 章会讲这台精密机器；堆肥桶发热的原理和酵母糖水一模一样：微生物在拆分子、释放能量，和酵母是同行。你在花园里施的每一把肥，都是在给植物体内看不见的反应补充原料。

\textbf{工业。}你早餐的面包（酵母的二氧化碳把它吹出蜂窝）、酸奶（乳酸菌把乳糖变成乳酸）、酱油和醋（漫长的微生物发酵）、味精、抗生素、味精之外的一大半现代药物——背后都是“选对微生物、控制反应条件、让正确的化学变化发生”的学问。第一册后面会学到，化工工程师的核心工作就是控制变化的方向和速度；而方向与速度由什么决定，是第 27 到 31 章的重头戏。

不过，在兴冲冲往前冲之前，有一个问题必须老实交代。这一整章，我们嘴上说“分子”，可你从来没有见过任何一个分子。我们说糖水里糖分子还在、焦糖里分子变了、气泡是二氧化碳——\textbf{凭什么？}看不见、摸不着的东西，怎么证明它存在，又怎么证明它是“这一种”而不是“那一种”？这不是抬杠，这是整个科学大厦的地基问题。下一章，我们就来做真正的侦探：学习怎样让看不见的东西自己留下证据。

\begin{tiaozhan}
“生成了新分子”——分子新在哪里？再往深处追一层：化学变化前后，每一种原子的“身份”都没变（碳还是碳，氧还是氧），改变的只有两样东西：原子之间的连接方式，和电子在原子之间的排布。前者叫化学键，是第三篇的主题；后者是氧化还原反应的实质，第 33 章揭晓。随之而来的问题是：原子凭什么肯重新组合？答案藏在两股力量的合谋里——新组合的能量是否更低，以及微观世界的概率规则是否允许（第 27、28 章）。还有一个让人失眠的数字游戏：仅用碳、氢、氧、氮四种原子，理论上能搭出的不同小分子，数量就远远超过全宇宙的原子总数。所以化学永远不可能靠“背完所有物质”来学会——你只能靠理解规则。这一段看不懂完全可以跳过，不影响主线；但如果它让你心里发痒，恭喜你，那是这本书最想保住的东西。
\end{tiaozhan}

\section*{练一练}

\begin{sikaoti}
\item 画三幅粒子图（用小圆圈代表分子）：①糖溶于水前后；②水烧开冒泡前后；③酵母糖水冒泡前后。用图说明：哪一幅里“分子种类”变了？图旁注明：圆圈只是表示法，真实的分子不是小球。
\item 给下面这些变化分类（物理变化/化学变化），并各写一句判据：冰雪融化、面包发霉、香水挥发、铁钉生锈、把面粉和水和成面团、面团被烤成面包。
\item 拧开汽水瓶立刻冒泡，酵母糖水慢慢冒泡。两个“冒泡”在粒子层面分别发生了什么？气泡里的分子各从哪里来？为什么说“冒泡”本身证明不了化学变化？
\item 数一数本章发酵反应式两边的碳、氢、氧原子个数，验证“账本”是否平衡。如果改成“一个葡萄糖分子只变成一个酒精分子”，账还能平吗？这说明反应式里的数字能随便改吗？
\item 一勺葡萄糖约含 $1.7\times 10^{22}$ 个分子。假如你每秒能数 100 个分子、昼夜不停，数完这一勺要多少年？（一年约 $3.2\times 10^{7}$ 秒。）算完回答：为什么化学家需要“摩尔”这种批发式计数法？
\item 在家里找三样东西（不许找危险物品），分别写下你观察到的变化线索，再写一句：这些线索可能怎样误导你？你打算用什么办法进一步确认？
\end{sikaoti}
